\section{The Rosenberg index} \label{sec:Rosenberg-index}
In this section, we provide a brief survey of the aspects of the Rosenberg index relevant to our statements and open problems.
More material on these matters can be found in the original articles of Rosenberg~\cite{Rosenberg:PSCNovikovI,Rosenberg:PSCNovikovII,Rosenberg:PSCNovikovIII} and in \cite{RosenbergStolz:PSCSurgeryConnections,Schick:ICM,Stolz:ManifoldsOfPSC,Stolz98Concordance}.

The starting point is that the spinor Dirac operator on a spin manifold is closely related to the scalar curvature via the Schr\"odinger--Lichnerowicz formula \(\Dirac^2 = \nabla^\ast \nabla + \frac{\scal}{4}\).
Together with the Atiyah--Singer index theorem, this was first used by Lichnerowicz~\cite{Lichnerwociz:Spineurs} to show that the \(\Ahat\)-genus is an obstruction to the existence of a positive scalar curvature metric on a closed spin manifold \(M\).
Later, this connection between the Dirac operator and scalar curvature was exploited further using variations of the index invariant.
For closed spin manifolds, the most sophisticated such invariant is the \emph{Rosenberg index} \(\alpha(M) \in \KO_n(\Cstar \pi_1 M)\) which takes values in the (real) K-theory of the group \textCstar-algebra \(\Cstar \pi_1 M\).
Notable examples of manifolds with non-vanishing Rosenberg index (but which in general do not have non-vanishing \(\Ahat\)-genus) include (area-)enlargeable spin manifolds~\cite{Hanke-Kotschick-Roe-Schick:Coarse-Top-Enlargeability, HankeSchick:Enlargeable,HankeSchick:EnlargeableInfinite}.

We will now describe two (equivalent) points of view towards the Rosenberg index.
The first is as the index of the spinor Dirac operator of \(M\) twisted by the ``universal flat bundle'' on~\(M\).
The latter is the \emph{Mi\v{s}\v{c}enko bundle} \(\mathcal{L}_M \to M\) which is the flat bundle of finitely generated projective Hilbert-\(\Cstar \Gamma\)-modules associated to the representation of \(\Gamma \coloneqq \pi_1 M\) on \(\Cstar \Gamma\) by left-multiplication.
Index theory extends to the situation of elliptic operators with coefficients in \textCstar-algebras yielding index invariants in corresponding operator K-theory groups.
This goes back to Mi\v{s}\v{c}enko and Fomenko~\cite{Mishchenko-Fomenko:Index}; for more recent expositions of the technical background relevant to our situation we refer to~\cite{Ebert:EllipticRegularityDirac,HankePapeSchick:CodimensionTwoIndex}.
The Rosenberg index \(\alpha(M) \in \KO_n(\Cstar \pi_1 M)\) is then simply the (K-theoretic) index of the twisted Dirac operator on \(\SpinBdl_M \otimes \mathcal{L}_M\) using the flat connection induced on the \Mishchenko\ bundle \(\mathcal{L}_M\), and \(\SpinBdl_M\) is the spinor bundle.
The Schr\"odinger--Lichnerowicz formula extends to this situation so that the non-vanishing of \(\alpha(M) \in \KO_n(\Cstar \Gamma)\) is still an obstruction to the existence of a positive scalar curvature metric on \(M\).
